\documentclass[11pt]{article}
\usepackage{amsmath,amssymb,amsthm,mathtools}

\newtheorem{definition}{Definition}
\newtheorem{proposition}{Proposition}
\newtheorem{lemma}{Lemma}
\newtheorem{corollary}{Corollary}
\newtheorem{remark}{Remark}

\newcommand{\RMain}{R^{\mathrm{red}}_{\beta,\Lambda,\mathrm{bg}}}
\newcommand{\MMain}{\mathsf M^{\mathrm{red}}_{\mathrm{Main}}}
\newcommand{\MOdd}{\mathsf M^{\mathrm{red}}_{\mathrm{odd}}}
\newcommand{\XiOne}{\Xi^{(1)}}
\newcommand{\rhobar}{\overline{\rho}^{\mathrm{dom}}}

\title{Basic Lemmas for Main-Visible Equivalence}
\date{2026-03-18}

\begin{document}
\maketitle

\section*{Purpose}

This note records the basic formal lemmas for the Main-visible quotient relation on a fixed
comparison stratum. It does not prove accepted-record factorization and does not assume any
minimality claim beyond the explicit readout map chosen here.

\begin{definition}[Fixed stratum]
Fix:
\begin{enumerate}
\item an accepted welded branch \(\beta\),
\item its accepted-tick/cutoff skeleton \(\{\Lambda_k\}_{k\in K_{\mathrm{acc}}(\beta)}\),
\item a background package \(\mathrm{bg}\) containing the common ambient weak Main-facing data
  \((E^-_{\Lambda_k},\mathcal A^-_{\Lambda_k},\mathsf B_{\Lambda_k})\), orientation convention, and
  any normalization convention used to pass between \(\rhobar_{\Lambda_k}\) and \(\chi_{\Lambda_k}\),
\item a common reduced welded boundary projector / reduced receiver target.
\end{enumerate}
Write \(\RMain\) for the set of all live admissible UOM seam/completion realizations \(\tau\) on this
fixed background.
\end{definition}

\begin{definition}[Reduced Main-visible readout]
For \(\tau\in\RMain\), define
\[
\MMain(\tau)
:=
\Big(
\mathcal P^{\mathrm{acc}}_{\Lambda_k}[\tau],\,
[\XiOne_{\Lambda_k}](\tau),\,\rhobar_{\Lambda_k}(\tau),\,\varepsilon_{\Lambda_k}(\tau)
\Big)_{k\in K_{\mathrm{acc}}(\beta)}.
\]
\end{definition}

\begin{definition}[Secondary odd-lane compression]
For \(\tau\in\RMain\), define
\[
\MOdd(\tau)
:=
\Big(
[\XiOne_{\Lambda_k}](\tau),\,\rhobar_{\Lambda_k}(\tau),\,\varepsilon_{\Lambda_k}(\tau)
\Big)_{k\in K_{\mathrm{acc}}(\beta)}.
\]
\end{definition}

\begin{definition}[Main-visible equivalence]
For \(\tau_1,\tau_2\in\RMain\), define
\[
\tau_1 \sim_{\mathrm{Main}} \tau_2
\qquad\Longleftrightarrow\qquad
\MMain(\tau_1)=\MMain(\tau_2).
\]
\end{definition}

\begin{proposition}[\(\sim_{\mathrm{Main}}\) is an equivalence relation]
The relation \(\sim_{\mathrm{Main}}\) is an equivalence relation on \(\RMain\).
\end{proposition}

\begin{proof}
\textbf{Reflexive.}
For every \(\tau\in\RMain\), one has \(\MMain(\tau)=\MMain(\tau)\), so
\(\tau\sim_{\mathrm{Main}}\tau\).

\textbf{Symmetric.}
If \(\tau_1\sim_{\mathrm{Main}}\tau_2\), then by definition
\(\MMain(\tau_1)=\MMain(\tau_2)\).
Equality is symmetric, so \(\MMain(\tau_2)=\MMain(\tau_1)\), hence
\(\tau_2\sim_{\mathrm{Main}}\tau_1\).

\textbf{Transitive.}
If \(\tau_1\sim_{\mathrm{Main}}\tau_2\) and \(\tau_2\sim_{\mathrm{Main}}\tau_3\), then
\[
\MMain(\tau_1)=\MMain(\tau_2)
\qquad\text{and}\qquad
\MMain(\tau_2)=\MMain(\tau_3).
\]
By transitivity of equality,
\(\MMain(\tau_1)=\MMain(\tau_3)\), so \(\tau_1\sim_{\mathrm{Main}}\tau_3\).
\end{proof}

\begin{lemma}[Equivalent \(\chi\)-presentation under fixed normalization]
Assume the normalization factors \(\nu_{\Lambda_k}\) are fixed in the background package
\(\mathrm{bg}\) and that
\[
\chi_{\Lambda_k}=\nu_{\Lambda_k}^{-1}\rho^{\mathrm{dom}}_{\Lambda_k}
\]
on the fixed stratum.
Then the relation defined by the tuple
\[
\Big(
\mathcal P^{\mathrm{acc}}_{\Lambda_k},\,
[\XiOne_{\Lambda_k}],\,\chi_{\Lambda_k},\,\varepsilon_{\Lambda_k}
\Big)_{k\in K_{\mathrm{acc}}(\beta)}
\]
coincides with \(\sim_{\mathrm{Main}}\).
\end{lemma}

\begin{proof}
On the fixed stratum, each \(\nu_{\Lambda_k}\) is constant.
Therefore equality of \(\rhobar_{\Lambda_k}\) is equivalent to equality of \(\chi_{\Lambda_k}\) at
each accepted tick.
The other components of the tuple are unchanged.
So the two readout maps have exactly the same fibers, hence define the same equivalence relation.
\end{proof}

\begin{corollary}[Accepted-record-history equality and \(\sim_{\mathrm{Main}}\) coincide]
On the fixed stratum, \(\tau_1\sim_{\mathrm{Main}}\tau_2\) if and only if \(\tau_1\) and \(\tau_2\)
have the same accepted-record history.
\end{corollary}

\begin{proof}
Equality of \(\MMain(\tau_1)\) and \(\MMain(\tau_2)\) includes equality of the first component
\(\mathcal P^{\mathrm{acc}}_{\Lambda_k}\) at every accepted tick, which is exactly equality of
accepted-record history.
Conversely, the accepted-record reduction theorem in the companion source note implies that equal
accepted-record histories force equality of every theorem-visible reduced receiver datum on the fixed
stratum, hence of the remaining components of \(\MMain\).
\end{proof}

\begin{corollary}[Downstream Main odd data descend to the quotient]
Assume the fixed background package contains the ambient odd spaces, ambiguity subspaces, and
Hilbertizing forms required by the polarized characterization interface.
Then any two realizations in the same \(\sim_{\mathrm{Main}}\)-class induce the same downstream odd
input for the finite-cutoff canonical representative theorem.
\end{corollary}

\begin{proof}
By definition of \(\sim_{\mathrm{Main}}\), realizations in the same class have the same descendant
class, the same descended scalar/orientation data, and the same orientation sign at every accepted
tick.
By definition of the fixed stratum, they also share the same ambient odd spaces, ambiguity subspaces,
and Hilbertizing forms.
These are exactly the odd-layer inputs used by the downstream Main theorem for the canonical
representative.
\end{proof}

\begin{remark}
The results above are purely formal.
They do not assert that \(\sim_{\mathrm{Main}}\) is the unique or minimal quotient beyond the fixed
readout chosen here.
The further compression from \(\MMain\) to \(\MOdd\) is a separate problem and is not settled by the
formal lemmas in this note.
\end{remark}

\end{document}
